%&pdflatex
\documentclass[11pt,a4paper]{article}

\usepackage[margin=1in,headheight=14pt]{geometry}
\usepackage[utf8]{inputenc}
\usepackage[T1]{fontenc}
\usepackage{microtype}
\usepackage{lmodern}
\usepackage{amsmath,amssymb}
\usepackage{booktabs}
\usepackage{array}
\usepackage{tabularx}
\usepackage{longtable}
\usepackage{enumitem}
\usepackage{fancyhdr}
\usepackage[dvipsnames,table]{xcolor}
\usepackage{listings}
\usepackage[colorlinks=true,linkcolor=MidnightBlue,urlcolor=NavyBlue,citecolor=ForestGreen]{hyperref}

% Document metadata
\newcommand{\toolname}{\texttt{vgrep}}
\newcommand{\toolversion}{v1.0.0}

% Header & Footer
\pagestyle{fancy}
\fancyhf{}
\fancyhead[L]{\textsc{\toolname\ Manual}}
\fancyhead[R]{\nouppercase{\leftmark}}
\fancyfoot[C]{\thepage}
\renewcommand{\headrulewidth}{0.4pt}
\renewcommand{\footrulewidth}{0.2pt}

% Section styling (pure standard LaTeX, no titlesec dependency required)
\makeatletter
\renewcommand{\section}{\@startsection{section}{1}{\z@}%
  {-3.5ex \@plus -1ex \@minus -.2ex}%
  {2.3ex \@plus .2ex}%
  {\normalfont\Large\bfseries\color{MidnightBlue}}}
\renewcommand{\subsection}{\@startsection{subsection}{2}{\z@}%
  {-3.25ex \@plus -1ex \@minus -.2ex}%
  {1.5ex \@plus .2ex}%
  {\normalfont\large\bfseries\color{RoyalBlue}}}
\renewcommand{\subsubsection}{\@startsection{subsubsection}{3}{\z@}%
  {-3.25ex \@plus -1ex \@minus -.2ex}%
  {1.5ex \@plus .2ex}%
  {\normalfont\normalsize\bfseries\color{CadetBlue}}}

% Ensure vertical mode across TOC generation to prevent horizontal-mode \addvspace errors
\renewcommand\tableofcontents{%
  \par
  \section*{\contentsname
    \@mkboth{\MakeUppercase\contentsname}{\MakeUppercase\contentsname}}%
  \par
  \@starttoc{toc}%
  \par
}

% Safeguard TOC entries so \addvspace always runs in vertical mode
\let\orig@l@section\l@section
\renewcommand*\l@section[2]{\par\orig@l@section{#1}{#2}}
\let\orig@l@subsection\l@subsection
\renewcommand*\l@subsection[2]{\par\orig@l@subsection{#1}{#2}}
\let\orig@l@subsubsection\l@subsubsection
\renewcommand*\l@subsubsection[2]{\par\orig@l@subsubsection{#1}{#2}}
\makeatother

% JSON language definition for listings
\lstdefinelanguage{json}{
    string=[s]{"}{"},
    stringstyle=\color{OliveGreen},
    comment=[l]{//},
    morecomment=[s]{/*}{*/},
    commentstyle=\color{gray}\itshape,
    literate=
     *{:}{{{\color{MidnightBlue}{:}}}}{1}
      {,}{{{\color{MidnightBlue}{,}}}}{1}
      {\{}{{{\color{RoyalBlue}{\{}}}}{1}
      {\}}{{{\color{RoyalBlue}{\}}}}}{1}
      {[}{{{\color{RoyalBlue}{[}}}}{1}
      {]}{{{\color{RoyalBlue}{]}}}}{1},
}

% Listings configuration
\lstdefinestyle{codeblock}{
    backgroundcolor=\color{black!4},
    basicstyle=\small\ttfamily,
    breaklines=true,
    breakatwhitespace=true,
    frame=single,
    rulecolor=\color{black!15},
    numbers=none,
    keywordstyle=\bfseries\color{RoyalBlue},
    commentstyle=\itshape\color{gray},
    stringstyle=\color{OliveGreen},
    showstringspaces=false,
    tabsize=4
}
\lstset{style=codeblock}

% Portable callout boxes (standard xcolor + minipage, no tcolorbox dependency required)
\newenvironment{callout}[2][MidnightBlue]{%
  \par\vspace{0.6em}\noindent
  \setlength{\fboxsep}{8pt}%
  \setlength{\fboxrule}{0.8pt}%
  \fcolorbox{#1!50}{#1!4}{%
    \begin{minipage}{\dimexpr\linewidth-2\fboxsep-2\fboxrule\relax}%
      {\bfseries\color{#1}#2}\par\vspace{0.4em}\small
}{%
    \end{minipage}%
  }\par\vspace{0.6em}%
}

\newenvironment{infobox}[1][Note]{\begin{callout}[MidnightBlue]{#1}}{\end{callout}}
\newenvironment{warnbox}[1][Warning]{\begin{callout}[RedOrange]{#1}}{\end{callout}}
\newenvironment{tipbox}[1][Pro Tip]{\begin{callout}[ForestGreen]{#1}}{\end{callout}}

% Keystroke macro
\newcommand{\keys}[1]{\fcolorbox{gray!40}{gray!10}{\small\texttt{\bfseries #1}}}

\begin{document}

\begin{titlepage}
    \centering
    \vspace*{1.5cm}
    
    {\Huge\bfseries\color{MidnightBlue} \toolname\ Manual\par}
    \vspace{0.4cm}
    {\Large\itshape Fast, Interactive, Project-Aware CLI Search and Refactoring Tool\par}
    \vspace{0.8cm}
    {\textsc{Version \toolversion}\par}
    
    \vspace{1.2cm}
    \begin{center}
    \setlength{\fboxsep}{12pt}%
    \setlength{\fboxrule}{0.8pt}%
    \fcolorbox{black!25}{black!3}{%
        \begin{minipage}{0.82\textwidth}
            \centering
            \texttt{\$ vgrep handleRequest\_fn}\\[0.25cm]
            \texttt{\$ vgrep --go "type .* struct"}\\[0.25cm]
            \texttt{\$ vgrep -v}
        \end{minipage}%
    }
    \end{center}
    
    \vspace{1.5cm}
    \begin{abstract}
        \noindent
        \toolname\ bridges terminal-based code search and high-speed editing. Built directly on top of ripgrep (\texttt{rg}), it circumvents shell quoting friction via an in-TUI raw input prompt, provides full Vim-motion navigation across search matches, supports granular in-place find-and-replace, and dispatches directly into editors with exact line and column coordinates.
    \end{abstract}
    
    \vfill
    {\small Compiled: \today\par}
\end{titlepage}

\tableofcontents
\newpage

\section{Introduction and Architecture}
\subsection{Core Philosophy}
\toolname\ is designed around four core operational principles:
\begin{itemize}[leftmargin=1.5em]
    \item \textbf{Frictionless Navigation}: Provide native Vim motions (\keys{j}, \keys{k}, \keys{J}, \keys{K}, \keys{g}, \keys{G}, counts, and relative numbers) for fluid movement across matches.
    \item \textbf{Direct Raw Input}: Eliminate shell quoting friction by reading keystrokes in raw terminal mode (\keys{n}), allowing arbitrary regex, spaces, and quotation marks without shell escaping.
    \item \textbf{Scoped Project Awareness}: Track search patterns automatically scoped to the active workspace or Git repository.
    \item \textbf{Exact Editor Positioning}: Seamlessly open editors with column-level precision and return directly to the search session upon exit.
\end{itemize}

\subsection{The Shell Escaping Dilemma}
Standard terminal pipelines involving \texttt{rg}, \texttt{grep}, or \texttt{find} suffer from intrinsic shell interpretation overhead:
\begin{itemize}[leftmargin=1.5em]
    \item \textbf{Eaten Quotations}: Shells (\texttt{bash}, \texttt{zsh}) strip or expand single (\texttt{'}) and double (\texttt{"}) quotes before passing arguments to executables. Searching for JSON snippets or string literals (e.g., \texttt{\{"key": "val"\}}) demands nested escaping.
    \item \textbf{Space Splitting}: Phrases with whitespace trigger parameter splitting unless strictly enclosed in quotes.
    \item \textbf{Expansion Collisions}: Metacharacters like \texttt{\$}, \texttt{!}, \texttt{\`{}}, and \texttt{\textbackslash} trigger history or variable expansions.
\end{itemize}

\toolname\ resolves this by reading keystrokes directly in terminal raw mode (\texttt{stty raw -echo}) inside its interactive Text User Interface (TUI). Characters entered during interactive search prompts (\keys{n}) are received verbatim without shell preprocessing.

\subsection{Architectural Overview}
The system architecture comprises five primary subsystems:
\begin{enumerate}[leftmargin=1.5em]
    \item \textbf{Ripgrep Execution Engine}: Subprocesses \texttt{rg} with \texttt{--json} and streams structured match records via standard output pipes.
    \item \textbf{Session and Quickfix Exporter}: Converts NDJSON streams into a clean JSON array written to \texttt{\textasciitilde/.config/wig/rg\_search.json} for cross-tool interoperability with editors such as \texttt{wig}.
    \item \textbf{Terminal Driver and Line Editor}: Handles alternate screen buffers (\texttt{\textbackslash033[?1049h}), raw terminal state restoration, SIGWINCH terminal resizing, and rune-level cursor manipulation.
    \item \textbf{Interactive Find \& Replace Engine}: Performs in-memory previewing and atomic file writes for batch refactoring.
    \item \textbf{Project-Aware History Store}: Automatically detects project roots (\texttt{.git}, \texttt{go.mod}, \texttt{Cargo.toml}, etc.) and manages localized query frequencies.
\end{enumerate}

\section{Installation and Dependencies}
\subsection{System Prerequisites}
\begin{table}[h!]
\centering
\small
\begin{tabularx}{\textwidth}{llX}
\toprule
\textbf{Dependency} & \textbf{Requirement} & \textbf{Functionality} \\
\midrule
\texttt{rg} (ripgrep) & \textbf{Mandatory} & Core regex search engine and file scanning \\
Go $\ge 1.25.0$ & Build only & Compiling from source code \\
\texttt{fzf} & Optional & Fuzzy picker for query history \\
\texttt{rgr} (repgrep) & Optional & External regex replace mode (\keys{r}) \\
\texttt{wig} / \texttt{nvim} / \texttt{vim} & Recommended & Text editor for match jumps and quickfix navigation \\
\bottomrule
\end{tabularx}
\caption{System dependencies and requirements}
\end{table}

\subsection{Building from Source}
Clone the repository and compile via the provided \texttt{Makefile}:
\begin{lstlisting}[language=bash]
git clone https://github.com/your-username/vgrep.git
cd vgrep
make build
sudo make install PREFIX=/usr/local
\end{lstlisting}

Verify the installation using the integrated health inspection command:
\begin{lstlisting}[language=bash]
vgrep --health
\end{lstlisting}

\section{Command Line Interface (CLI)}
\subsection{Command Syntax}
\begin{lstlisting}[language=bash]
vgrep [FLAGS] [SEARCH_PATTERN]
\end{lstlisting}

\subsection{Command Line Flags}
\begin{table}[h!]
\centering
\small
\begin{tabularx}{\textwidth}{llX}
\toprule
\textbf{Flag} & \textbf{Alias} & \textbf{Description} \\
\midrule
\texttt{-i} & \texttt{--ignore-case} & Execute case-insensitive search. \\
\texttt{-v} & \texttt{--view} & View matches from the previous session without re-scanning. \\
\texttt{-e} & \texttt{--edit} & Open \texttt{\textasciitilde/.config/vgrep/config.toml} in configured editor. \\
\texttt{--health} & & Inspect dependencies, tool versions, and configuration paths. \\
\texttt{--init} & & Purge history file and cached search session. \\
\texttt{--fzf} & & Force fallback to the \texttt{fzf} interface instead of the full TUI. \\
\midrule
\texttt{--go} & & Limit search to Go source files (\texttt{*.go}). \\
\texttt{--rs} & & Limit search to Rust source files (\texttt{*.rs}). \\
\texttt{--py} & & Limit search to Python source files (\texttt{*.py}). \\
\texttt{--cc} & & Limit search to C/C++ source files (\texttt{*.c}, \texttt{*.cpp}, \texttt{*.h}, etc.). \\
\texttt{--dart} & & Limit search to Dart source files (\texttt{*.dart}). \\
\texttt{--swift} & & Limit search to Swift source files (\texttt{*.swift}). \\
\bottomrule
\end{tabularx}
\caption{Command line flag reference}
\end{table}

\subsection{Function Shorthand Expansion (\texttt{\_fn})}
Appending the \texttt{\_fn} suffix to any search query triggers automatic language-aware keyword conversion based on detected project root artifacts:
\begin{itemize}[leftmargin=1.5em]
    \item If \texttt{go.mod} is discovered or \texttt{--go} is supplied:
    \begin{center}
    \texttt{vgrep handleRequest\_fn} $\;\longrightarrow\;$ \texttt{"func handleRequest"}
    \end{center}
    \item In Rust projects or generic workspaces:
    \begin{center}
    \texttt{vgrep handleRequest\_fn} $\;\longrightarrow\;$ \texttt{"fn handleRequest"}
    \end{center}
\end{itemize}

\subsection{Single Match Auto-Jump}
When a search query yields \textbf{exactly one match} across the entire project, \toolname\ skips TUI presentation and directly opens the configured editor at the exact line and column of the match.

\section{Interactive TUI Reference}
When multiple matches are identified, \toolname\ launches an alternate screen TUI.

\subsection{Screen Layout}
The interface is divided into four functional zones:
\begin{enumerate}[leftmargin=1.5em]
    \item \textbf{Title Bar}: Shows active query, flags (\texttt{[-i]}, \texttt{[-F]}), language filters, filter substring, and aggregate match/file counters.
    \item \textbf{Viewport}: Scrollable list of matched files (headers) and individual match lines with relative line numbering and syntax highlights.
    \item \textbf{Status Line}: Mode indicator (\texttt{NORMAL}, \texttt{SEARCH}, \texttt{FILTER}, \texttt{REPLACE}, \texttt{PREVIEW}), count buffers, selected file/line/column, and viewport position.
    \item \textbf{Command \& Help Bar}: Key bindings and dynamic interactive prompts.
\end{enumerate}

\subsection{Normal Mode Keybindings}
\begin{table}[h!]
\centering
\small
\begin{tabularx}{\textwidth}{lX}
\toprule
\textbf{Key} & \textbf{Action} \\
\midrule
\keys{j} / \keys{k} & Move selection cursor down / up by 1 entry. \\
\keys{<count>j} / \keys{<count>k} & Move selection cursor down / up by \texttt{<count>} entries (e.g., \keys{5j}). \\
\keys{J} / \keys{l} & Jump forward to the first match of the next file. \\
\keys{K} / \keys{h} / \keys{L} & Jump backward to the first match of the previous file. \\
\keys{g} / \keys{<count>g} & Jump to the first entry (row 0) or to absolute row \texttt{<count>}. \\
\keys{G} / \keys{<count>G} & Jump to the last entry or to absolute row \texttt{<count>}. \\
\keys{Ctrl+D} / \keys{Ctrl+F} & Page down by viewport height. \\
\keys{Ctrl+U} / \keys{Ctrl+B} & Page up by viewport height. \\
\keys{Enter} / \keys{o} / \keys{e} & Open selected entry in editor at line and column. \\
\keys{SPC} (Space) & Toggle inclusion/exclusion of current match or entire file. \\
\keys{a} & Toggle inclusion/exclusion for \textbf{all} matches across all files. \\
\keys{/} & Enter incremental in-memory live filter mode. \\
\keys{c} & Clear active live filter. \\
\keys{n} & Open interactive prompt for a new ripgrep search query. \\
\keys{F} & Toggle between literal string search (\texttt{-F}) and regular expression search. \\
\keys{Alt+i} & Toggle case-sensitivity (\texttt{-i}) on the fly. \\
\keys{Tab} / \keys{R} & Enter interactive in-app Find \& Replace mode. \\
\keys{r} & Launch external \texttt{rgr} (repgrep) interactive replacer (if installed). \\
\keys{q} / \keys{Ctrl+C} & Cleanly restore terminal and exit. \\
\bottomrule
\end{tabularx}
\caption{Normal mode navigation and command matrix}
\end{table}

\subsection{Search Prompt Mode (\keys{n})}
Pressing \keys{n} in normal mode presents the search prompt:
\begin{itemize}[leftmargin=1.5em]
    \item Type search patterns freely without shell quotation or space escapes.
    \item \keys{Alt+i}: Toggle case-sensitivity flag \texttt{[-i]} dynamically.
    \item \keys{Enter}: Run ripgrep, record query to project history, and refresh TUI.
    \item \keys{Esc} / \keys{Ctrl+C}: Dismiss prompt without altering active results.
\end{itemize}

\subsection{Live Filter Mode (\keys{/})}
Filter mode narrows displayed results in-memory without invoking \texttt{rg}:
\begin{itemize}[leftmargin=1.5em]
    \item Substring filtering matches against both file paths and line contents.
    \item \keys{\(\uparrow\)} / \keys{\(\downarrow\)}: Navigate matches while keeping the filter prompt open.
    \item \keys{Ctrl+U}: Wipe current filter input.
    \item \keys{Enter} / \keys{Esc}: Confirm filtered view and return to Normal mode.
\end{itemize}

\section{Find and Replace Workflow}
\subsection{Interactive In-Place Replacement (\keys{Tab} / \keys{R})}
\toolname\ features a native refactoring engine:
\begin{enumerate}[leftmargin=1.5em]
    \item \textbf{Trigger}: Press \keys{Tab} or \keys{R} in Normal mode.
    \item \textbf{Input}: Type the replacement string in the \texttt{REPLACE>} prompt.
    \item \textbf{Inspect (\keys{Tab})}: Switch to \texttt{PREVIEW} mode. Matching text highlights with substitution previews in green (\textcolor{ForestGreen}{\texttt{HighlightSubst}}).
    \item \textbf{Exclusion}: Use \keys{SPC} on individual matches or headers (or \keys{a} for all) to mark entries with strikethrough red badges (\texttt{-[REMOVED]}). Excluded lines remain untouched.
    \item \textbf{Commit (\keys{Enter})}: Writes changes directly to disk preserving original file line-endings (CRLF or LF) and file permissions. Updated results are immediately refreshed in the TUI.
\end{enumerate}

\subsection{External Replacer Integration (\keys{r})}
When \texttt{rgr} (repgrep) is installed in \texttt{\$PATH}, pressing \keys{r} temporarily suspends the TUI and launches \texttt{rgr} with current pattern and file type filters. When \texttt{rgr} exits, \toolname\ resumes seamlessly.

\section{Editor Integration}
\subsection{Command Formatting by Editor}
\toolname\ inspects the executable name of your configured editor and formats command-line arguments to position the cursor at the exact match line and column:

\begin{table}[h!]
\centering
\small
\begin{tabularx}{\textwidth}{llX}
\toprule
\textbf{Editor} & \textbf{Command Pattern} & \textbf{Format Example} \\
\midrule
\texttt{wig} & \texttt{wig <file> +<line>:<col>} & \texttt{wig main.go +42:16} \\
\texttt{nvim} / \texttt{vim} / \texttt{vi} & \texttt{vim "+call cursor(<line>, <col>)" <file>} & \texttt{nvim "+call cursor(42, 16)" main.go} \\
\texttt{code} (VS Code) & \texttt{code -g <file>:<line>:<col>} & \texttt{code -g main.go:42:16} \\
\texttt{hx} / \texttt{helix} & \texttt{hx <file>:<line>:<col>} & \texttt{hx main.go:42:16} \\
\texttt{nano} & \texttt{nano +<line>,<col> <file>} & \texttt{nano +42,16 main.go} \\
\texttt{emacs} & \texttt{emacs +<line>:<col> <file>} & \texttt{emacs +42:16 main.go} \\
Generic Fallback & \texttt{<editor> +<line> <file>} & \texttt{subl +42 main.go} \\
\bottomrule
\end{tabularx}
\caption{Editor column positioning protocol}
\end{table}

\subsection{Editor Resolution Hierarchy}
The active editor is resolved via the following sequence:
\begin{enumerate}[leftmargin=1.5em]
    \item The \texttt{editor} key in \texttt{\textasciitilde/.config/vgrep/config.toml} (if executable exists in \texttt{\$PATH}).
    \item The system \texttt{\$EDITOR} environment variable (if executable exists).
    \item Auto-probed binaries in priority: \texttt{wig} $\longrightarrow$ \texttt{nvim} $\longrightarrow$ \texttt{vim}.
    \item Fallback default: \texttt{"vim"}.
\end{enumerate}

\subsection{Shared Session File Format (\texttt{rg\_search.json})}
Every search writes a structured array of match records to the session file (default \texttt{\textasciitilde/.config/wig/rg\_search.json}):
\begin{lstlisting}[language=json]
[
  {
    "file_path": "/home/user/project/main.go",
    "line": 42,
    "char": 15,
    "text": "\tkeys := wig.NewKeyHandler()\n"
  }
]
\end{lstlisting}

This shared file can be parsed by editors or loaded back into \toolname\ via:
\begin{lstlisting}[language=bash]
vgrep -v
\end{lstlisting}

\section{Configuration Specification}
The configuration file is located at \texttt{\textasciitilde/.config/vgrep/config.toml}. It can be edited directly using \texttt{vgrep -e}.

\begin{lstlisting}[language=bash]
# ~/.config/vgrep/config.toml

# Preferred editor executable name or absolute path
editor = "wig"

# Target path for quickfix session JSON export
session_file = "~/.config/wig/rg_search.json"

# Enable literal matching (-F) by default
fixed_strings = false
\end{lstlisting}

\begin{table}[h!]
\centering
\small
\begin{tabularx}{\textwidth}{llX}
\toprule
\textbf{Key} & \textbf{Type} & \textbf{Semantics} \\
\midrule
\texttt{editor} & string & Executable name or path for opening files. \\
\texttt{session\_file} & string & File path for quickfix session export (supports \texttt{\textasciitilde} expansion). \\
\texttt{fixed\_strings} & boolean & When \texttt{true}, searches use literal mode (\texttt{-F}) by default. \\
\bottomrule
\end{tabularx}
\caption{Configuration options reference}
\end{table}

\section{Project-Scoped History}
When invoked without query arguments (\texttt{vgrep}), the tool scans parent directories for marker files (\texttt{.git}, \texttt{go.mod}, \texttt{Cargo.toml}, \texttt{pubspec.yaml}, \texttt{pyproject.toml}, \texttt{package.json}) to determine the project root.

\begin{itemize}[leftmargin=1.5em]
    \item Queries are stored in \texttt{\textasciitilde/.config/vgrep/history.json} partitioned by project root path.
    \item Up to 25 queries are tracked per repository, ranked by frequency and relative recency (e.g., \textit{just now}, \textit{10m ago}, \textit{3d ago}).
    \item If \texttt{fzf} is installed, an interactive fuzzy picker is displayed. Otherwise, a numbered terminal menu is provided.
    \item Item \texttt{0} allows direct resumption of the previous session (equivalent to \texttt{vgrep -v}).
\end{itemize}

\section{Diagnostic Health Check}
The \texttt{--health} flag inspects installed tooling, active configuration, and filesystem paths:

\begin{lstlisting}[language=bash]
$ vgrep --health
=== vgrep Health Check ===

  ripgrep (rg): ripgrep 14.1.0 (~/.cargo/bin/rg)
  fzf: Found (/usr/bin/fzf)
  rgr (repgrep): Found (~/.cargo/bin/rgr)
  Config editor (wig): Found (/usr/local/bin/wig)

  Config path:  ~/.config/vgrep/config.toml
  History path: ~/.config/vgrep/history.json
  Wig session:  ~/.config/wig/rg_search.json
\end{lstlisting}

\end{document}